\subsection{A positive model for the containing iterated tensor}

The obstruction in \fid{F63} concerns a specified simple
summand. Its containing four-dimensional tensor admits
an explicit positive model with the actual named action.
The simple projector and a single hard-Lefschetz center
still impose stronger requirements.

Put
\[
 \delta=v+v^{-1},\qquad D=\delta^2,\qquad
 a=D-2,\qquad d=D^2,\qquad A=D^2-8D+8,
\]
and let \(R=\Q[v,v^{-1}]_{(v-1)}\).
Retain the actual restriction
\(N(D-2)\otimes N(D-2)\) through
\(\Delta^{(2,2)}\) from \fid{F63}.
The iteration parameter is \(k=4\); the strand rank
remains \(r=3\).
In the left-column convention its generators are
\begin{equation}\label{eq:positiveIteratedActualMatrices}
 R_1=\begin{pmatrix}
 d&Da&Da&2a^2\\0&0&0&-Da\\0&0&0&-Da\\0&0&0&d
 \end{pmatrix},\qquad
 R_2=\begin{pmatrix}
 d&0&0&0\\-Da&0&0&0\\-Da&0&0&0\\2a^2&Da&Da&d
 \end{pmatrix}.
\end{equation}
These are the source restriction matrices, with the
normalization verified from Definition~2.5,
equations~(8)--(11), p.~6, equation~(13), p.~8,
and Theorems~3.4 and~3.7, pp.~12--14
of \cite{BlasiakChebyshev}.

\begin{finding}[\fid{F66}, a regular positive whole-tensor model]
\label{f:iteratedPositiveWholeTensor}
The actual four-dimensional module
\eqref{eq:positiveIteratedActualMatrices} has an
\(R\)-basis, fixed by its transported tensor bar,
in which both named generators have entries in
\(\N[v^2,v^{-2}]\).
It admits concrete graded matrix kernels with coherent
quadratic comparison maps and a pure weight-zero
ordinary-Tate derived realization.

The original split Grothendieck image has generic
dimension six. After full Karoubi closure, that image
is the full matrix algebra of generic dimension
sixteen, and the original kernel category excludes
the transpose-dual adjoint.
The specified simple projector is not effective.
Moreover the actual unmarked two-word trace has no
degree-two single-center Lefschetz realization.
\end{finding}

\begin{proof}
The proof below gives the basis, categorical maps,
closure computations and arithmetic realization
explicitly. The Lefschetz claim is proved separately
after the construction.

\par\smallskip\noindent\textbf{The positive actual basis.}\quad
Set
\begin{equation}\label{eq:positiveIteratedXY}
 x=(D-2)^2=v^4+2+v^{-4},\qquad
 y=4(D-1)=4v^2+4+4v^{-2}.
\end{equation}
Thus \(x+y=d\) and \(x-y=A\).
Consider
\begin{equation}\label{eq:positiveIteratedMatrices}
 S_1=\begin{pmatrix}
 x&y&0&0\\x&y&0&0\\0&0&y&x\\0&0&y&x
 \end{pmatrix},\qquad
 S_2=\begin{pmatrix}
 x&0&y&0\\0&y&0&x\\x&0&y&0\\0&y&0&x
 \end{pmatrix}.
\end{equation}
All entries are effective even Laurent polynomials.
Let \(C\) have columns
\[
 c_+=(1,1,1,1)^T,\quad c_-=(y,-x,-x,y)^T,
 \quad f_1=(1,1,-1,-1)^T,\quad f_2=(1,-1,1,-1)^T.
\]
Multiplication gives the complete action table
\[
\begin{array}{c|cccc}
 &c_+&c_-&f_1&f_2\\\hline
 S_1&dc_+&0&df_1&Af_1\\
 S_2&dc_+&0&Af_2&df_2.
\end{array}
\]
For example,
\(S_1f_2=(x-y,x-y,y-x,y-x)^T=Af_1\).
The determinant is \(\det C=-8(x+y)=-8d\), a unit
at \(v=1\).

Let \(B\) have the columns already verified in
\fid{F63}:
\[
 (d,-Da,-Da,d)^T,\quad(0,1,-1,0)^T,\quad
 (A,-Da,-Da,d)^T,\quad(d,-Da,-Da,A)^T.
\]
Direct multiplication gives
\[
 B^{-1}R_iB=\operatorname{diag}(d,0,N_i(A)),\qquad
 N_1(A)=\begin{pmatrix}d&A\\0&0\end{pmatrix},\quad
 N_2(A)=\begin{pmatrix}0&0\\A&d\end{pmatrix},
\]
and \(\det B=128D(D-2)(D-1)^2\), also a unit
at \(v=1\).
The action table therefore proves
\begin{equation}\label{eq:positiveIteratedIntertwiner}
 T=BC^{-1}\in\GL_4(R),\qquad T^{-1}R_iT=S_i.
\end{equation}
This retains the named action and the same local
lattice, not only the module dimension.

The original tensor basis is bar-fixed by the binary
basis construction in \fid{F43}/\fid{F56}.
Every entry of \(B,C\), and hence \(T\), is a
rational expression in \(D\), fixed by
\(v\mapsto v^{-1}\). The new basis is therefore
bar-fixed as well. This is the transported formal bar;
it is not a positive Hermitian polarization or a
claim about VIII's global integral lattice.

\par\smallskip\noindent\textbf{The actual projector.}\quad
In this positive generator basis the projector onto
\(N(A)\) is
\begin{equation}\label{eq:positiveIteratedProjector}
 E=C\operatorname{diag}(0,0,1,1)C^{-1}
  =\frac12\begin{pmatrix}
   1&0&0&-1\\0&1&-1&0\\0&-1&1&0\\-1&0&0&1
  \end{pmatrix}.
\end{equation}
It kills \(c_+,c_-\), fixes \(f_1,f_2\), and
commutes with both \(S_i\). In particular
\[
                         \Tr(ES_1S_2)=A^2.
\]
The mixed Laurent coefficients proved in \fid{F63}
exclude a basis making \(E,S_1,S_2\) all effective
Laurent matrices: their product and trace would then
have only nonnegative coefficients.
The displayed constant projector exhibits the
negative entries explicitly.
An arbitrary linear direct summand does not thereby
become an effective object of the coefficient-kernel
category.

\par\smallskip\noindent\textbf{Concrete kernels and quadratic maps.}\quad
Let \(\mathcal V_{\mathrm{even}}\) be finite graded
vector spaces in even internal degrees, with
degree-zero maps, over a characteristic-zero field.
Choose \(X,Y\) with characters \(x,y\) in
\eqref{eq:positiveIteratedXY}:
\(X\) has one line in each degree \(4,-4\) and
two in degree zero; \(Y\) has four lines in each
degree \(2,0,-2\).
Set \(Z=X\oplus Y\), so \(\operatorname{ch}Z=d\).
Replace \(x,y,0\) in
\eqref{eq:positiveIteratedMatrices} by \(X,Y,0\)
to define kernels \(K_1,K_2\).
An entry \(K_{ab}\) sends state \(b\) to state
\(a\), and convolution is
\begin{equation}\label{eq:positiveIteratedConvolution}
                   (K\circ L)_{ac}
                         =\bigoplus_b K_{ab}\otimes L_{bc}.
\end{equation}
It matches the left-column matrix convention.
The associators and pentagon are the actual
finite-sum/tensor identifications before taking
characters.

For \(K_1\), the blocks are \(\{1,2\}\) and
\(\{3,4\}\); the rows agree within each block.
For \(K_2\), the blocks are \(\{1,3\}\) and
\(\{2,4\}\), again with equal rows.
The two coefficient entries are \(X,Y\), in one
order or the other. Summing the intermediate state
therefore supplies actual distributivity isomorphisms
\begin{equation}\label{eq:positiveIteratedQuadratic}
                             K_i^2\simeq Z\otimes K_i.
\end{equation}
The coefficient is exactly the iterated quadratic
coefficient \(d=\delta^4\).

For a length-\(r\) path using only \(K_i\), label
the \(r-1\) intermediate states in path order.
Their independent choices give the explicit map
\begin{equation}\label{eq:positiveIteratedPowers}
                       K_i^r\simeq
                          Z^{\otimes(r-1)}\otimes K_i.
\end{equation}
Raw quadratic reductions may extract the \(Z\)
factors in different orders. Already at length three,
the two reductions differ by their transposition.
Compare each reduction with
\eqref{eq:positiveIteratedPowers} by the canonical
coefficient permutation restoring the labeled path
order, together with associativity.
These comparisons agree: adjacent transpositions
satisfy the symmetric-group relations, and all maps
send the same labeled path factors to the same
ordered positions. This uses ordinary symmetric
coherence; it does not assert identity of the two
uncorrected length-three maps.

Mixed words already compose in the concrete matrix
category. Their characters give the actual \(S_i\)
products, so all source algebra relations hold after
taking characters by
\eqref{eq:positiveIteratedIntertwiner}.
No source-prescribed higher categorical presentation,
or categorical iteration functor, follows from these
quadratic comparisons alone.

\par\smallskip\noindent\textbf{Split \(K_0\), completion and adjoints.}\quad
Let \(\mathcal P\) be the full additive monoidal
subcategory of \(\operatorname{Mat}_4(\mathcal
V_{\mathrm{even}})\) generated by the identity,
\(K_1,K_2\) and even shifts, with all ambient graded
kernel morphisms.
Two objects are isomorphic exactly when all graded
entry dimensions agree. Its additive isomorphism
monoid thus embeds in the cancellative monoid
\(\N[v^2,v^{-2}]^{16}\).
If two formal differences have equal character,
adding the opposite terms gives objects with the
same character, hence an actual isomorphism.
Consequently
\[
 K_0^{\mathrm{split}}(\mathcal P)
       \lhook\joinrel\longrightarrow
            \operatorname{Mat}_4(\Z[v^2,v^{-2}])
\]
is injective. Its image is exactly the algebra
generated by \(I,S_1,S_2\), since its objects are
finite sums of shifted words.

Generically that algebra is
\(\Q(v)\times\Q(v)\times M_2(\Q(v))\).
Indeed the projector is in the generated algebra:
\[
 E=\frac{d(S_1+S_2)-S_1S_2-S_2S_1}{d^2-A^2}.
\]
The plus projector is \((I-E)S_1/d\), and its
complement within \(I-E\) is the minus projector.
The two-dimensional block is absolutely irreducible
because \(A\ne0\) and \(A^2\ne d^2\).
To see its full matrix algebra directly, diagonalize
the rank-one idempotent \(N_1/d\).
Both off-diagonal entries of \(N_2/d\) in that
basis are nonzero under these two inequalities.
Multiplication by the two diagonal idempotents
isolates the off-diagonal matrix units; together
they span \(M_2\).
The generic dimension is therefore six.

In the full Karoubi closure the identity splits into
the four state units \(E_{aa}\).
The kernel \(E_{aa}K_iE_{bb}\) isolates an entry.
Every nonzero entry \(X\) or \(Y\) contains a
degree-zero line, which splits by an ambient graded
idempotent. Hence the closure contains \(E_{ab}\)
for every directed support edge.
The supports include both directions of
\(1\)--\(2\), \(3\)--\(4\),
\(1\)--\(3\), and \(2\)--\(4\).
Composing along paths yields every matrix unit.
Even shifts and sums then yield all of
\(\operatorname{Mat}_4(\mathcal V_{\mathrm{even}})\),
already idempotent complete.
Its split group has rank sixteen over
\(\Z[v^2,v^{-2}]\), rather than six generically.

The transpose-graded-dual of \(K_1\) is absent
from \(\mathcal P\). Every matrix in the original
algebra preserves the common line spanned by
\(c_+\), since \(S_ic_+=dc_+\).
The coefficients are bar-invariant, so the adjoint
character is \(S_1^T\), but
\[
                       S_1^Tc_+=(2x,2y,2y,2x)^T
\]
is not proportional to \(c_+\), since \(x-y=A\ne0\).
This character is outside the original image.
Adjoints exist in the completed ambient category;
the assertion concerns the original action image.

\par\smallskip\noindent\textbf{Ordinary-Tate pure derived kernels.}\quad
Fix \(\F_Q\), \(\ell\nmid Q\), and a finite
coefficient field \(E/\Q_\ell\).
Give multiplicity spaces trivial Frobenius and put
\begin{equation}\label{eq:positiveIteratedTateFunctor}
                \Phi(V)=\bigoplus_{e\in2\Z}
                    V_e\otimes E(-e/2)[-e].
\end{equation}
Each summand has its degree-\(e\) cohomology pure
of point weight \(e\), so it is a pure complex of
weight zero by BBD~5.1.8, pp.~126--127
\cite{BBD}.
All shifts are even. Degree addition and tensor
symmetry therefore make \(\Phi\) symmetric strong
monoidal, without a Koszul sign correction.
For every extension degree,
\[
       \Tr_{\mathrm{Eul}}(F^n;\Phi(V))
                       =\operatorname{ch}(V)(Q^{n/2}).
\]
This uses only ordinary integral Tate twists;
there is no odd half-Tate motive.
Apply it entrywise to
\eqref{eq:positiveIteratedConvolution}.
Finite sums and tensor products preserve point
purity of weight zero, so the actual quadratic
and repeated-power maps lift to pure derived
kernels. Every closed-word trace object is pure
of weight zero with the actual unmarked word trace.

Entrywise graded duality maps to ordinary point
duality, since the dual of \(E(-e/2)[-e]\) is
\(E(e/2)[e]\). Symmetric self-dual choices for
\(X,Y\) retain the formal bar. This is not the
transpose-dual adjoint excluded above.
\end{proof}

\paragraph{The operation and source boundaries.}
Purity of complex weight zero is not one common
Frobenius radius on total cohomology: the degree-\(e\)
group has point weight \(e\).
No path-sign bound, dual spectral bound, designated
canonical labels or effective simple projector
is supplied by this construction.
It is a coefficient-object realization of changing
compositions, not a fixed finite Weil space carrying
all word traces.

The regular basis change of the actual restriction
does not itself define a monoidal functor between
the original binary kernel categories and this
four-state category. The source formula involves
\(Q_i=D-P_i\), virtual in the strict effective
kernel category. The positive replacement of the
whole restricted module does not make those
\(Q_i\) effective or give an all-\(k\) comparison.
The tested module occurs in the 4096-dimensional
tensor cube of the fourfold ordinary two-dimensional
tensor product,
\[
                   (V_1\otimes V_2\otimes V_3\otimes V_4)^{\otimes3},
                       \qquad \dim V_i=2.
\]
No positive model for all its other summands is
claimed. The full Karoubi closure changes the
algebra image and introduces individual state
objects; it does not identify those with the
source's canonical simple labels.

\subsubsection{The whole-tensor Lefschetz obstruction}

The preceding construction has pure point
complexes, but a single degree-two Lefschetz
center would impose an additional condition.
The actual unmarked two-word trace is
\begin{align}\label{eq:positiveIteratedWholeTrace}
 T(v)&=\Tr(S_1S_2)=d^2+A^2\\
 &=2v^8+40v^4+64v^2+108
                  +64v^{-2}+40v^{-4}+2v^{-8}.\notag
\end{align}
This follows directly from the scalar-plus,
scalar-minus and \(N(A)\) blocks, and the
expansion has value \(320\) at \(v=1\).
It is effective, but its coefficients in
degrees \(6\) and \(-6\) vanish.

\begin{finding}[\fid{F66}, the additional single-center test]
\label{f:iteratedPositiveLefschetzGate}
No graded vector space with dimension polynomial
\(T(v)\), or any Laurent monomial shift of it,
admits a degree-two Lefschetz operator about one
center. In particular \(T\) is not a genuine
\(\mathrm{SL}_2\) character, so both actual
generators cannot be realized by matrix kernels
of genuine finite-dimensional
\(\mathrm{SL}_2\) representations with the
stated composition/closed-trace comparison.

There is no projective structure map
\(f:X\to\Spec\F_Q\) of a finite-type scheme
and pure perverse sheaf \(P\) of integral weight
such that \(Rf_*P\) has Euler traces
\(T(Q^{n/2})\) at every sufficiently large
extension degree.
\end{finding}

\begin{proof}
A degree-two Lefschetz operator centered at an
integer \(c\) requires
\[
                   L^i:V^{c-i}\xrightarrow{\sim}V^{c+i}
                         \qquad(i\ge0).
\]
For \(T\), its support extrema \(-8,8\)
force \(c=0\) by this symmetry.
But \(\dim V^{-8}=2\) and \(V^{-6}=0\).
The first step of \(L^8:V^{-8}\to V^8\)
is zero, a contradiction.
For \(v^sT\), the extrema force \(c=s\),
and the identical gap argument applies.

For \(\mathrm{SL}_2\), the highest visible
weight is eight with multiplicity two.
Each highest-weight-eight constituent contains
a weight-six line; larger highest weights
cannot occur without visible top terms.
The missing weight six is impossible.
Equivalently the standard raising operator
would give the contradicted Lefschetz maps.
If the generator entries were genuine
\(\mathrm{SL}_2\) representations, direct
sums, tensor products and the sum of diagonal
entries would give a genuine representation
with character \(T\). This proves the
kernel assertion without a simple projector.

For the arithmetic assertion, first let \(K\)
be any pure point complex of integral weight
\(w\) with the stated eventual traces.
Finite power-sum uniqueness follows by taking
sufficiently large consecutive integers and
using the Vandermonde determinant.
The visible roots are \(Q^{e/2}\), of point
weight \(e\), with multiplicities the
coefficients of \(T\).
BBD~5.1.8 puts each root only in degree
\(e-w\) \cite{BBD}.
Thus a fixed root cannot cancel between
degrees, and hidden roots cannot cancel either.
All visible coefficients are positive and all
\(e\) are even, so \(w\) is even and
\[
                 \dim H^{e-w}(K)=[v^e]T(v),
\]
with no other cohomology. This allows Jordan
blocks of arithmetic Frobenius.

For \(K=Rf_*P\), proper direct image
preserves purity by \(f_!=f_*\) and dual
stability, as in BBD~5.1.13--14, p.~128
\cite{BBD}.
An ample class supplies the relative
hard-Lefschetz isomorphisms
\[
          L^i:H^{-i}(K)\xrightarrow{\sim}H^i(K)(i),
                               \qquad i\ge0,
\]
by BBD, Theorem~5.4.10, p.~144 \cite{BBD}.
The support extrema \(-8-w,8-w\) must be
symmetric about zero, so \(w=0\).
Then \(H^{-8}\) has dimension two and
\(H^{-6}=0\), contradicting the first step
of \(L^8\).
\end{proof}

The support-gap proof also excludes a nonzero
positive constant multiple whenever it has
integral coefficients, and a Laurent monomial
grading shift. Arbitrary Laurent polynomial
normalization is not covered.
The projective statement requires a pure
\emph{perverse} source and one Lefschetz center.
It does not exclude arbitrary pure complexes
on proper schemes: the explicit pure point
complex above, with the identity map of a
point, already demonstrates that boundary.
Separately shifted projective pieces with
different centers are not excluded.
An independently inserted endomorphism trace
also need not have positive multiplicities.
None of these extra Lefschetz comparisons is
asserted to be a published canonical-basis axiom.

% Separate paragraph for Part III.
\paragraph{The positive tensor and its Lefschetz boundary in Grothendieck's language.}
The regular intertwiner and named-action comparison
are finite locally free module data.
EGA~I, Theorem~1.3.7, p.~86, Corollary~1.3.8,
p.~87 and Theorem~1.4.1, pp.~90--92 provide
the affine module dictionary \cite{EGA}.
SGA~3, I, 4.7.2--3 and cote~90, scan~17 supply
action language \cite{SGA3,Archive}.
The even torus grading is weaker than an
\(\mathrm{SL}_2\) action with a raising operator.

SGA~6, V, \S\S2--3 and cote~133, scan~18
supply coefficient representation and
\(\lambda/\sigma\) language for the symmetric
coefficient category, not a \(\lambda\)-ring
structure on the noncommutative word algebra
\cite{SGA6,Archive}.
The difference between the six-dimensional action
image and sixteen-dimensional completed matrix
algebra is the explicit closure calculation above.

Ordinary Tate lines and bounded complexes fit
SGA~6, I, Definition~4.7, pp.~121--122 and
Example~4.8, p.~122 \cite{SGA6}.
BBD~5.1.5 and~5.1.8, pp.~125--127 fix the
pure-complex convention \cite{BBD}.
The stronger projective-perverse conclusion
uses BBD~5.1.13--14, p.~128 and
Theorem~5.4.10, p.~144.
These sources do not identify the auxiliary
states with canonical simples or supply the
additional single-center comparison.
The ordinary pure coefficient construction
and its Lefschetz limitation are both proved
here with their distinct hypotheses.
